% !TeX root = ../main.tex
\chapter{System Requirements}
\label{chap:system-requirements}
\lhead{\emph{System Requirements}}

\section{Functional Requirements}

The system ingests alerts from external SIEM sources, normalises them to a common schema, enriches them with contextual data, and produces ranked incident candidates. For each incident the scoring component generates an explainable risk score and confidence value, accompanied by reason codes and context multipliers. The planning component then produces a structured response plan---strategy, priority, action set, and rollback steps---which is validated against explicit policy rules and compliance-oriented checks before any action is dispatched. The execution stage records action outcomes and errors.

Analyst feedback is a first-class input. An analyst can mark an incident as true positive, false positive, or needs more information, and optionally override severity. Feedback events are recorded in an append-only log and used to update model weights in a bounded, auditable way. Investigation traces are tracked and used to recommend next steps. A user interface exposes raw alerts, normalised alerts, scoring explanations, plans, policy checks, execution status, and audit logs.

Figure~\ref{fig:req-traceability} provides a requirement-to-component traceability view,
showing where each requirement class is enforced in the architecture.

\begin{figure}[H]
  \centering
  \includegraphics[width=\textwidth]{Diagrams/Requirements_Traceability.png}
  \caption{Traceability mapping from functional, non-functional, and safety requirements to core system components.}
  \label{fig:req-traceability}
\end{figure}

\section{Non-Functional Requirements}

The architecture is modular so that intelligence models can be swapped without touching orchestration code. The interface between layers is explicit and versioned. The system must run as a single project to facilitate demonstrations and repeated experiments. Explanations are human-readable and consistent across alerts to allow direct comparison.

There is a clear separation between domain logic, application services, infrastructure, and the UI. Structured logging throughout facilitates audits and retrospective analysis.

\section{Security and Safety Requirements}

Automation is bounded by design. All planner actions are subject to policy rules---allow/deny lists, confidence thresholds, asset criticality limits. High-impact actions require approval gates, and each automated action has a corresponding rollback. The execution layer validates parameters, enforces entity-presence checks, and refuses to execute when required context is missing or policy checks fail. Rate limits and idempotency safeguards prevent repeated execution under noisy alert conditions.

If confidence is too low or policy approval is not granted, the system falls back to recommendation-only behaviour. Model risk is controlled through bounded, logged weight updates. Model output is traceable after updates, and the system supports drift detection and periodic offline recalibration. Connectors and executors operate under least-privilege access, and sensitive fields are minimised or redacted where possible. All policy evaluations and execution outcomes are logged with timestamps and actor identity.
